% Source PDF pages 89--92; manual pages 2-60--2-63.
\subsection{Gravity Model}\label{sec:gravity-model}

The gravitational field surrounding a planet can be represented by a potential
function $U$. The acceleration required for trajectory propagation is the gradient
of that potential:
\begin{equation}
  \nabla U=\vec a_{\mathrm{grav}}.
  \label{eq:gravity-potential-gradient}
\end{equation}
For the earth, $U$ is the geopotential. It is a function of the geocentric coordinates
$\lVert\vec r\rVert$, $\delta_{gc}$, and $\lambda$ and is commonly expanded in
spherical harmonics as\manualref{6}
\begin{equation}
  \begin{aligned}
  U=\frac{GM}{\lVert\vec r\rVert}
  \Bigg\{1
  &+\sum_{n=2}^{\infty}\sum_{m=0}^{n}
    \left[\frac{R_{\oplus}}{\lVert\vec r\rVert}\right]^n
    P_{n,m}(\sin\delta_{gc})\\[-0.2em]
  &\qquad\times
    \left[C_{n,m}\cos(m\lambda)+S_{n,m}\sin(m\lambda)\right]
  \Bigg\}.
  \end{aligned}
  \label{eq:geopotential-spherical-harmonics}
\end{equation}
The coefficients $C_{n,m}$ and $S_{n,m}$ are related to spatial features of the
earth's gravity field. Earth gravity models derived from satellite tracking therefore
consist of values for $GM$, $R_{\oplus}$, $C_{n,m}$, and $S_{n,m}$.

Because the geocentric origin coincides with the earth's center of mass, the
first-degree terms $C_{1,0}$ and $C_{1,1}$ are zero. The expansion therefore begins
at degree $n=2$.

\taossubhead{Geopotential}

The geopotential can be interpreted as three parts. The first is the point-mass term
\begin{equation}
  \frac{GM}{\lVert\vec r\rVert},
  \label{eq:point-mass-geopotential}
\end{equation}
which is the complete potential used by the spherical-earth model. All harmonic
corrections are neglected for that model.

The second part contains terms that depend on latitude but not longitude, namely the
$m=0$ terms. These are the zonal harmonics. The degree-2 zonal harmonic, which
represents the earth's oblateness, is the largest spherical-harmonic correction. A
zonal-only potential is
\begin{equation}
  U=\frac{GM}{\lVert\vec r\rVert}
  \left\{1+\sum_{n=2}^{\infty}
    \left[\frac{R_{\oplus}}{\lVert\vec r\rVert}\right]^n
    P_{n,0}(\sin\delta_{gc})C_{n,0}\right\}.
  \label{eq:zonal-geopotential-c-coefficients}
\end{equation}
Defining $J_n=-C_{n,0}$ gives the more common form
\begin{equation}
  U=\frac{GM}{\lVert\vec r\rVert}
  \left\{1-\sum_{n=2}^{\infty}
    \left[\frac{R_{\oplus}}{\lVert\vec r\rVert}\right]^n
    P_{n,0}(\sin\delta_{gc})J_n\right\}.
  \label{eq:zonal-geopotential-j-coefficients}
\end{equation}

The third part consists of the remaining longitude-dependent terms. The largest is
usually the degree-2, order-2 term, which represents the earth's asymmetry about the
equator. Even that term is substantially smaller than the degree-2 zonal harmonic
associated with oblateness.\manualref{10}

\taossubhead{Legendre Functions}

The terms $P_{n,m}(\sin\delta_{gc})$ are Legendre functions of
$\sin\delta_{gc}$ defined by\manualref{6}
\begin{equation}
  P_{n,m}(\sin\delta_{gc})
  =\cos^m\delta_{gc}\,
   \frac{d^m}{d(\sin\delta_{gc})^m}
   \left[P_{n,0}(\sin\delta_{gc})\right],
  \label{eq:associated-legendre-definition}
\end{equation}
where
\begin{equation}
  P_{n,0}(\sin\delta_{gc})
  =\frac{1}{2^n n!}
   \frac{d^n}{d(\sin\delta_{gc})^n}
   \left(\sin^2\delta_{gc}-1\right)^n.
  \label{eq:legendre-order-zero-definition}
\end{equation}
The notation
\begin{equation}
  \frac{d^n}{d(\sin\delta_{gc})^n}
  \left[P(\sin\delta_{gc})\right]
  \label{eq:legendre-derivative-notation}
\end{equation}
denotes the $n$th derivative of $P$ with respect to $\sin\delta_{gc}$. The first
four order-zero functions are
\begin{equation}
  \begin{aligned}
  P_{1,0}(\sin\delta_{gc})&=\sin\delta_{gc},\\
  P_{2,0}(\sin\delta_{gc})&=\tfrac12
    \left(3\sin^2\delta_{gc}-1\right),\\
  P_{3,0}(\sin\delta_{gc})&=\tfrac12
    \left(5\sin^3\delta_{gc}-3\sin\delta_{gc}\right),\\
  P_{4,0}(\sin\delta_{gc})&=\tfrac18
    \left(35\sin^4\delta_{gc}-30\sin^2\delta_{gc}+3\right).
  \end{aligned}
  \label{eq:first-four-legendre-functions}
\end{equation}

\taossubhead{Normalization}

The spherical-harmonic coefficients in
\cref{eq:geopotential-spherical-harmonics,eq:zonal-geopotential-c-coefficients}
are unnormalized. As degree increases, the unnormalized coefficients become very
small partly because of the physical gravity field and partly because the associated
Legendre functions become large. A common normalization scales the functions by
\begin{equation}
  \overline{P}_{n,m}
  =\left[(2-\delta_m)(2n+1)
    \frac{(n-m)!}{(n+m)!}\right]^{1/2}P_{n,m},
  \label{eq:normalized-legendre-function}
\end{equation}
and scales a cosine coefficient by the inverse factor:
\begin{equation}
  \overline{C}_{n,m}
  =\left[\frac{1}{(2-\delta_m)(2n+1)}
    \frac{(n+m)!}{(n-m)!}\right]^{1/2}C_{n,m}.
  \label{eq:normalized-gravity-coefficient}
\end{equation}
Here $\delta_m=1$ for $m=0$ and $\delta_m=0$ for $m>0$. A bar denotes a
normalized function or coefficient.

TAOS uses unnormalized Legendre functions and coefficients because it retains no
terms above degree 4. Published coefficients for the WGS-84\manualref{6} and
GEM-T1\manualref{11} models are normalized and must be unnormalized before use
in TAOS.

\taossubhead{Gravity Acceleration Vector}

Because $\vec a_{\mathrm{grav}}=\nabla U$, its components in geocentric coordinates
are
\begin{equation}
  \vec a_{\mathrm{grav}}\mathbin{\cdot}\hat x_{gc}
  =\frac{1}{\lVert\vec r\rVert}
   \frac{\partial U}{\partial\delta_{gc}},
  \label{eq:gravity-geocentric-x-gradient}
\end{equation}
\begin{equation}
  \vec a_{\mathrm{grav}}\mathbin{\cdot}\hat y_{gc}
  =\frac{1}{\lVert\vec r\rVert\cos\delta_{gc}}
   \frac{\partial U}{\partial\lambda},
  \label{eq:gravity-geocentric-y-gradient}
\end{equation}
and
\begin{equation}
  \vec a_{\mathrm{grav}}\mathbin{\cdot}\hat z_{gc}
  =-\frac{\partial U}{\partial\lVert\vec r\rVert}.
  \label{eq:gravity-geocentric-z-gradient}
\end{equation}
The radial component is negated because the position vector $\vec r$ points in the
opposite direction from the inward geocentric unit vector $\hat z_{gc}$.
Differentiating \cref{eq:geopotential-spherical-harmonics} gives
\begin{equation}
  \begin{aligned}
  \vec a_{\mathrm{grav}}\mathbin{\cdot}\hat x_{gc}
  ={}&\frac{GM}{\lVert\vec r\rVert^2}
  \sum_{n=2}^{\infty}
  \left[\frac{R_{\oplus}}{\lVert\vec r\rVert}\right]^n
  \sum_{m=0}^{n}
  \frac{\partial P_{n,m}(\sin\delta_{gc})}{\partial\delta_{gc}}\\[-0.2em]
  &\qquad\times
  \left[C_{n,m}\cos(m\lambda)+S_{n,m}\sin(m\lambda)\right],
  \end{aligned}
  \label{eq:gravity-full-geocentric-x}
\end{equation}
\begin{equation}
  \begin{aligned}
  \vec a_{\mathrm{grav}}\mathbin{\cdot}\hat y_{gc}
  ={}&\frac{GM}{\lVert\vec r\rVert^2}
  \sum_{n=2}^{\infty}
  \left[\frac{R_{\oplus}}{\lVert\vec r\rVert}\right]^n
  \sum_{m=0}^{n}
  \frac{mP_{n,m}(\sin\delta_{gc})}{\cos\delta_{gc}}\\[-0.2em]
  &\qquad\times
  \left[-C_{n,m}\sin(m\lambda)+S_{n,m}\cos(m\lambda)\right],
  \end{aligned}
  \label{eq:gravity-full-geocentric-y}
\end{equation}
and
\begin{equation}
  \begin{aligned}
  \vec a_{\mathrm{grav}}\mathbin{\cdot}\hat z_{gc}
  =\frac{GM}{\lVert\vec r\rVert^2}
  \Bigg\{1
  &+\sum_{n=2}^{\infty}(n+1)
    \left[\frac{R_{\oplus}}{\lVert\vec r\rVert}\right]^n\\[-0.2em]
  &\quad\times\sum_{m=0}^{n}P_{n,m}(\sin\delta_{gc})
    \left[C_{n,m}\cos(m\lambda)+S_{n,m}\sin(m\lambda)\right]
  \Bigg\}.
  \end{aligned}
  \label{eq:gravity-full-geocentric-z}
\end{equation}

The function \texttt{accel\_grav} evaluates these expressions through degree
$n=4$ for the \texttt{wgs-84-full} and \texttt{gem-t1-full} earth models described
in \cref{sec:block-earth}. Solar and lunar gravity are neglected. The
\texttt{tsap-72} and \texttt{tsap-84} models retain zonal harmonics through degree
4 for compatibility with TSAP. Those models are somewhat inconsistent because the
$J_3$ and $J_4$ terms are roughly the same order of magnitude as some omitted
nonzonal terms.

The standard \texttt{wgs-72} and \texttt{wgs-84} models retain only the dominant
second-degree zonal harmonic $J_2$. Their acceleration equations reduce to
\begin{equation}
  \vec a_{\mathrm{grav}}\mathbin{\cdot}\hat x_{gc}
  =-\frac{GM}{\lVert\vec r\rVert^2}
    \left[\frac{R_{\oplus}}{\lVert\vec r\rVert}\right]^2
    \left(3J_2\sin\delta_{gc}\cos\delta_{gc}\right),
  \label{eq:gravity-j2-geocentric-x}
\end{equation}
\begin{equation}
  \vec a_{\mathrm{grav}}\mathbin{\cdot}\hat y_{gc}=0,
  \label{eq:gravity-j2-geocentric-y}
\end{equation}
and
\begin{equation}
  \vec a_{\mathrm{grav}}\mathbin{\cdot}\hat z_{gc}
  =\frac{GM}{\lVert\vec r\rVert^2}
   \left\{1-
    \left[\frac{R_{\oplus}}{\lVert\vec r\rVert}\right]^2
    \frac{3}{2}J_2\left(3\sin^2\delta_{gc}-1\right)
   \right\}.
  \label{eq:gravity-j2-geocentric-z}
\end{equation}
The $J_2$-only models provide adequate accuracy for a point-mass trajectory code
such as TAOS. For most trajectories, the effects of higher-order harmonics are small
relative to uncertainties in other models, especially the atmosphere.
